\section{总结与展望}
\label{sec:conclusion}

\subsection{论文主要工作}

本文以燃煤机组汽轮机热力系统为研究对象，围绕设备健康状态建模、故障知识图谱构建以及KG-GCN故障诊断方法开展研究。针对汽轮机系统在变工况运行下健康状态持续迁移、故障知识分散以及高质量故障样本有限等问题，构建了由动态状态残差和静态领域知识共同驱动的故障诊断流程。主要研究工作和结论如下。

（1）针对汽轮机设备正常运行状态随负荷、主蒸汽边界和操作调节持续变化的问题，采用设备级$U+B\rightarrow Y$建模方式，以实际DCS边界条件和操作变量为输入、设备状态参数为输出建立健康状态模型。选取VHP作为代表性对象，对ANN、TCN、iTransformer和GRU进行比较。2024年5月训练、2024年6月独立测试结果表明，GRU在最终24个状态参数上的宏平均$R^2$达到0.9511，整体性能优于其他对比模型。健康状态模型能够在给定运行边界下提供动态健康参考，实测状态与健康预测之间的偏差可进一步作为故障诊断的动态状态信息。

（2）针对汽轮机系统设备关系复杂以及故障知识分散于DCS资料、设备结构信息和故障文本中的问题，构建基于运行生产资料的汽轮机系统故障知识图谱。图谱以设备、参数、异常、故障和处置为主要实体类型，将机组DCS测点、设备IO和系统接口作为设备结构知识，并利用公开论文、标准和运行事件资料补充通用故障机理。通过大语言模型生成候选三元组，再结合实体类型、关系方向、设备兼容性、来源证据和Embedding语义对齐完成知识校核与实体标准化。当前知识库包含127个标准实体和123条候选关系，其中121条关系进入主知识图谱，另有2条存在资料差异的候选关系保留待核。由此实现了汽轮机设备结构、参数关联和故障机理的统一图结构表达。

（3）针对不同汽轮机故障共享局部监测参数、少量故障标签条件下分类模型容易产生误诊的问题，构建基于KG-GCN的故障诊断模型。该方法将健康状态模型得到的多参数残差窗口映射到知识图谱参数节点，通过两层GCN沿设备和参数关系逐层聚合特征，并在训练阶段首先忽略类别标签进行随机残差遮蔽重构；完成预训练后固定图卷积编码器，利用带标签故障窗口训练分类层。以VHP为代表性设备，基于真实健康运行残差和独立故障机理构造正常状态、VHP通流部件侵蚀、一级抽汽支路泄漏和VHP汽封泄漏四类半物理场景。固定随机种子主实验中，KG-GCN的Accuracy和F1分别达到98.57\%和98.58\%，高于CNN的97.86\%和97.85\%以及AE的95.54\%和95.54\%；3组随机种子重复实验中，KG-GCN平均Macro-F1为98.46\%$\pm$0.11\%。当分类阶段仅使用20\%带标签故障窗口时，掩码预训练使Macro-F1由94.48\%提高至96.67\%，表明该训练方式在故障标签有限条件下能够获得更加稳定的结构化残差表示。

综上，本文形成了“设备健康状态建模--状态残差构建--汽轮机领域知识图谱--GCN结构化特征聚合--掩码预训练--故障分类”的诊断流程。其中，知识图谱和KG-GCN方法面向汽轮机热力系统建立，VHP作为代表性设备用于当前算例验证。对于HP、IP、LP、凝汽器及其他回热设备，本文工作提供了相应的参数子图构建和模型扩展基础，但其故障诊断性能仍需通过独立设备数据和故障样本进一步验证。

\subsection{后续研究方向}

本文研究仍存在进一步完善空间。首先，当前故障分类验证采用真实健康残差背景和独立故障机理构造的半物理故障场景，尚不能替代现场真实故障事件验证。后续应结合机组缺陷记录、检修记录和历史故障数据，对故障发生时刻、故障严重程度以及多参数响应特征进行进一步校准，并在条件允许时利用经实测数据校准的热力仿真模型补充不同故障程度下的动态样本。

其次，当前算例仅完成VHP任务子图的故障诊断实验。虽然系统级知识图谱已经覆盖主通流、再热、抽汽回热和凝汽冷端等主要对象，但HP、IP、LP和凝汽器等设备尚未完成同等条件下的独立故障分类验证。后续可按照相同的参数节点映射和KG-GCN训练流程扩展至其他设备，并进一步研究跨设备故障传播、故障耦合及复合故障识别问题。

最后，汽轮机故障知识图谱需要随运行生产资料的积累持续更新。后续可进一步引入检修报告、运行日志和专家经验作为候选知识来源，在实体标准化、来源追溯和领域规则校核后完成增量更新；同时研究关系置信度、边权重及动态图结构对GCN诊断性能的影响，提高模型在复杂变工况、故障样本有限和监测信息不完整条件下的工程适应性。
